% ================================================================
% Section: SLMs as the Substrate of Agentic AI
% New section (position core of the consolidated paper)
% ================================================================

\section{Small Language Models as the Substrate of Agentic AI}
\label{sec:agentic}

The preceding sections establish three facts about the sub-10B regime: modern SLMs reach capability levels that were state-of-the-art for far larger models only two years earlier (Section~\ref{sec:reported}); they can be compressed to run on commodity and edge hardware with modest accuracy loss (Sections~\ref{sec:compression}--\ref{sec:deployment}); and a mature ecosystem now exists to deploy them (Section~\ref{subsec:deployment_frameworks}). This section argues that these facts, taken together with how agentic systems are actually built, make SLMs the \emph{default} execution substrate for AI agents, with large models reserved for a shrinking set of genuinely hard calls. The argument follows and extends the position of \citet{belcak2025slm_agentic}; we situate it within the broader literature on LLM--SLM collaboration \citep{chen2024role_small,chen2025collab_survey} and connect it to the deployment evidence consolidated in this paper.

\subsection{What Agents Delegate to Language Models}
\label{subsec:agentic_delegation}

An AI agent is a software system in which one or more language models operate under a control loop that also includes tools, memory, and a scaffold of prompts and parsers \citep{belcak2025slm_agentic}. The language model in such a system is rarely asked to hold an open-ended conversation. It is asked to perform bounded sub-tasks: decide which tool to call and with what arguments, extract fields from a document, classify an intent, rewrite a query, summarize a retrieved passage, produce a short code patch, or check whether a step succeeded. These calls share three properties. They are \emph{narrow}---each exercises a small slice of the model's competence. They are \emph{repetitive}---the same handful of call types recur thousands of times in production. And they are \emph{format-constrained}---the output must parse as JSON, match a function signature, or fall within a fixed label set, or the surrounding code breaks \citep{belcak2025slm_agentic}.

None of these properties rewards the breadth of a frontier LLM. A model that can also write a sonnet, debate philosophy, and reason across a 200-page contract is carrying capacity that the agent never invokes on the call in question. \citet{belcak2025slm_agentic} formalize this as the observation that agents use only a narrow subset of language-model functionality, and that agentic systems demand strict behavioral alignment to code interfaces---exactly the setting in which a specialized small model can match or exceed a general large one.

\subsection{The Case for SLMs in Agents}
\label{subsec:agentic_case}

\subsubsection{Capability Sufficiency}
\label{subsubsec:agentic_capability}

The first premise is that SLMs are already good enough for delegated sub-tasks. On instruction following, tool selection, and structured extraction---the operations that dominate agent traces---sub-10B models trained or tuned for the purpose are competitive with frontier models. Salesforce's xLAM family \citep{zhang2024xlam} illustrates the point: its 7B function-calling model ranked at the top of the Berkeley Function-Calling Leaderboard, and its 1B variant was reported to surpass GPT-3.5-Turbo and Claude-3~Haiku on the same task despite being one to two orders of magnitude smaller. Unified-reasoning SLMs extend this to multi-step tasks: Qwen3 \citep{qwen3_2025} exposes a controllable ``thinking budget'' so that an agent can spend inference compute only on the calls that need it, and distilled reasoning models such as Phi-4-mini-reasoning \citep{xu2025phi4mini_reasoning} and the DeepSeek-R1 distillations \citep{deepseek2025r1} bring competition-mathematics performance into the 3--8B range. \citet{belcak2025slm_agentic} argue that objections based on scaling laws (AV1: larger is permanently better) understate SLM capability because scaling-law studies hold architecture fixed, exclude task-specific fine-tuning, and predate inference-time reasoning---all three of which disproportionately help small models.

\subsubsection{Economics}
\label{subsubsec:agentic_economics}

The second premise is cost. \citet{belcak2025slm_agentic} estimate that serving a 7B SLM is 10--30$\times$ cheaper than serving a 70--175B LLM once latency, energy, and FLOPs are accounted for, and that parameter-efficient fine-tuning lets a small model be specialized for a new agent role in GPU-hours rather than GPU-weeks. The deployment figures consolidated in this paper are consistent with that range: Table~\ref{tab:latency_compare} shows a 5--20$\times$ cost advantage and comparable or better latency for a quantized 3B model relative to a 70B model on chatbot and code-completion workloads, and an absolute advantage in offline settings where the large model cannot run at all. Because an agent amortizes model cost over many calls per user task, a constant-factor saving on each call compounds into the difference between a viable and an unviable product margin.

\subsubsection{Behavioral Alignment and Operational Fit}
\label{subsubsec:agentic_alignment}

The third premise is control. An agent needs its model to emit exactly the format the scaffold expects, and to do so reliably across a long tail of inputs. Constraining a small specialized model to a narrow output distribution is easier than suppressing the much larger behavioral range of a frontier model, and fine-tuning a small model on the agent's own logged traces closes the residual gap \citep{belcak2025slm_agentic}. \citet{belcak2025slm_agentic} note a virtuous cycle here (their argument A7): agent interactions generate high-quality, in-distribution training data essentially for free, which is precisely what specialization of a small model requires. Operationally, small models are also easier to version, roll back, run on-premises for data-residency reasons, and co-locate with the tools they call.

\subsection{Heterogeneous SLM--LLM Systems}
\label{subsec:agentic_hetero}

The position is not that LLMs disappear from agents but that the default inverts: SLMs handle the repetitive language errands, and a strong generalist model is invoked only for planning, open-ended decomposition, or calls that genuinely need long context or broad world knowledge \citep{belcak2025slm_agentic}. Figure~\ref{fig:hetero_agent} sketches the arrangement. Realizing it requires a mechanism to decide, per call, which model to use. \citet{chen2025collab_survey} organize the collaboration literature into five interaction patterns---pipeline, routing, auxiliary, distillation, and fusion---of which routing is the most direct fit for agents.

\begin{figure}[t]
\centering
\footnotesize
\begin{tikzpicture}[
  >={Stealth[length=2mm]},
  box/.style={draw, rounded corners=2pt, align=center, inner sep=4pt, line width=0.5pt},
  slm/.style={box, fill=black!4, text width=40mm, align=left},
  llm/.style={box, fill=black!10, text width=40mm, align=left},
  router/.style={box, fill=black!6, text width=24mm, minimum height=10mm},
  lbl/.style={font=\scriptsize\itshape, text=black!55},
]
\node[box, text width=13mm] (task) {User task};
\node[box, right=7mm of task, minimum height=12mm, text width=20mm] (loop) {Agent\\control loop};
\node[box, below=6mm of loop, font=\scriptsize, text width=20mm] (tools) {Tools $\cdot$ Memory};
\node[router, right=9mm of loop] (router) {\textbf{Router}\\[1pt]\scriptsize difficulty $\cdot$ cost $\cdot$ privacy $\cdot$ network};
\node[slm, above right=2mm and 12mm of router] (slm) {\textbf{Specialised SLMs} — default\\[1pt]\scriptsize tool selection, argument filling, field extraction, classification, query rewrite, short code patch, step verification};
\node[llm, below right=2mm and 12mm of router] (llm) {\textbf{Generalist LLM} — escalation\\[1pt]\scriptsize planning, open-ended decomposition, long context, broad world knowledge};
\draw[->] (task) -- (loop);
\draw[->] (tools) -- (loop);
\draw[<->] (loop) -- (router);
\draw[->] (router.east) -- ($(slm.west)+(0,-1mm)$) node[lbl, midway, above] {most calls};
\draw[->] (router.east) -- ($(llm.west)+(0,1mm)$) node[lbl, midway, below] {minority};
\end{tikzpicture}
\caption{Heterogeneous SLM--LLM agent. A per-call router sends the majority of
delegated sub-tasks to specialised small models and escalates only planning-like
or long-context calls to a generalist LLM; tools and memory are shared. After
\citet{belcak2025slm_agentic}.}
\label{fig:hetero_agent}
\end{figure}

\subsubsection{Learned Routing}
\label{subsubsec:agentic_routing}

RouteLLM \citep{ong2024routellm} trains a lightweight router on human-preference data to predict, for each query, whether a weak model's answer would be judged as good as a strong model's, and dispatches accordingly; the authors report cost reductions of more than 2$\times$ on standard benchmarks with no measurable quality loss, and transfer of the trained router to unseen strong/weak model pairs. Division-of-Thoughts \citep{division_of_thoughts_2025} pushes routing below the query level, decomposing a task into sub-questions and assigning each to a local SLM or a cloud LLM according to estimated difficulty, for on-device agents. The open problems are difficulty estimation without an oracle, routing objectives that jointly weigh accuracy, latency, privacy sensitivity, and network availability, and handoff that does not produce a visible seam in output quality \citep{demystifying2025slm}.

\subsubsection{Production Instances}
\label{subsubsec:agentic_production}

Two large deployments already embody the heterogeneous pattern. Apple's Foundation Models framework pairs an on-device model of roughly 3B parameters, quantized with 2-bit quantization-aware training and KV-cache sharing, with a Private Cloud Compute server model built on a Parallel-Track Mixture-of-Experts transformer, and routes to the server only when the on-device model is expected to fall short \citep{gunter2025afm}. The on-device model is reported to be competitive with Qwen3-4B and Gemma-3-4B on its target tasks. Third-party providers can now register cloud models behind the same API surface as the local model, making the SLM the default and the LLM an opt-in escalation \citep{apple2026afm3}. Samsung's Gauss and Google's Gemini Nano occupy the same design point on their respective platforms \citep{samsung2023gauss}.

\subsection{Converting an LLM Agent to an SLM Agent}
\label{subsec:agentic_conversion}

For teams with an existing LLM-only agent, \citet{belcak2025slm_agentic} propose a six-step conversion pipeline that we summarize here because it operationalizes the position. \textbf{S1, instrumentation:} log every non-interface model call (inputs, outputs, latency) from the running agent, with encryption and access control. \textbf{S2, curation:} strip personally identifiable and protected health information and application secrets from the logs. \textbf{S3, clustering:} apply unsupervised clustering to the curated calls to surface the recurring request types---these are the specialization targets. \textbf{S4, selection:} for each cluster, choose a candidate SLM by matching required capability against benchmarks, license, and deployment constraints. \textbf{S5, specialization:} fine-tune the chosen SLM on the cluster's curated data, typically with LoRA or QLoRA. \textbf{S6, iteration:} re-run the loop periodically as the traffic distribution drifts. The pipeline is attractive precisely because the data it needs is a by-product of running the agent, and because each cluster converted removes a recurring frontier-model bill.

\subsection{Barriers to Adoption}
\label{subsec:agentic_barriers}

\citet{belcak2025slm_agentic} identify three reasons the shift has lagged the technical case. First, large centralized-inference investments create sunk-cost inertia and an operational bias toward one big model behind one API (B1). Second, SLMs are still evaluated on generalist leaderboards (MMLU, broad chat benchmarks) that are misaligned with agentic utility, so their fitness for delegated sub-tasks is under-measured (B2)---an evaluation gap this paper's consolidated tables also reflect, since public reporting concentrates on the same generalist benchmarks. Third, SLM releases receive far less attention than frontier launches, so awareness lags capability (B3). None of these is a technical obstacle, which is the point: the barriers are those of infrastructure, measurement, and attention, and each is already eroding as routing frameworks mature, agent-specific benchmarks such as the Berkeley Function-Calling Leaderboard gain traction, and on-device platforms ship SLMs to billions of devices by default.
